% 出席確認クイズ 第04回　企業活動と情報システムの全体像
%   attendance/attendance04.html から生成（解説は本ガイド用に加筆）。

\quizq{1}{ERP(統合基幹業務システム)の説明として適切なものはどれでしょうか?}%
  {顧客との関係を管理するシステム}%
  {\correct{会計・人事・在庫などの基幹業務を統合して管理するシステム}}%
  {電子メールを配信するシステム}%
  {社内向けのSNS}%
  {正解は (b)。ERP（Enterprise Resource Planning）は、会計・人事・生産・在庫などの基幹業務のデータを1つのデータベースで統合管理するシステムです。(a) は CRM の説明です。部門ごとにばらばらだったデータが統合されることで、全社の状況を一貫した数字で把握できるようになります。}

\quizq{2}{CRMが主に扱う対象はどれでしょうか?}%
  {原材料の調達}%
  {\correct{顧客との関係(購買履歴や問い合わせなど)}}%
  {建物の管理}%
  {従業員の給与}%
  {正解は (b)。CRM（Customer Relationship Management）は、顧客情報・購買履歴・問い合わせなどを管理し、顧客との関係を深めるためのしくみです。(a) の調達は SCM、(d) の給与は ERP の人事給与モジュールが担います。}

\quizq{3}{SCMの目的として適切なものはどれでしょうか?}%
  {社内文書を保管する}%
  {\correct{調達から生産・物流・販売までの供給の流れを最適化する}}%
  {ソフトウェアを開発する}%
  {顧客の満足度を調査する}%
  {正解は (b)。SCM（Supply Chain Management）は、調達から生産・物流・販売に至る供給の流れ全体を、在庫や納期の面で最適化します。(a)(c)(d) は文書管理・開発・調査で、供給連鎖の最適化とは別の活動です。}

\quizq{4}{「バリューチェーン(価値連鎖)」の考え方を提唱したのは誰でしょうか?}%
  {スティーブ・ジョブズ}%
  {\correct{マイケル・ポーター}}%
  {アダム・スミス}%
  {ピーター・ドラッカー}%
  {正解は (b)。バリューチェーン（価値連鎖）はマイケル・ポーターが『競争優位の戦略』（1985年）で示した枠組みで、企業活動を5つの主活動と4つの支援活動に分けて、どこで価値が生まれるかを分析します。ドラッカーはマネジメント論、スミスは分業論で知られ、ジョブズは経営者です。}

\quizq{5}{クラウド(SaaS)型の情報システムにAI機能が組み込まれることの利点として適切なものはどれでしょうか?}%
  {従業員の教育が不要になる}%
  {\correct{自社で開発しなくても最新のAI機能を利用できる}}%
  {データが社外に出ることが絶対にない}%
  {ネットワークが不要になる}%
  {正解は (b)。SaaS 型のサービスでは提供側が機能を更新するため、自社で開発しなくても最新のAI機能を使えます。一方で (c) は誤りで、データは事業者のサーバーに置かれるため契約と設定の確認が必要です。(a) 教育は引き続き必要で、(d) ネットワークはむしろ不可欠です。}
